SCRATCH TEXT.

MAYBE-SWAP

* [AGDA](https://parfunc.slack.com/files/vikraman/F5YGV9JQ1/Untitled.hs)
* [COQ](https://www.cs.princeton.edu/~appel/vfa/Perm.html#Exploration1.maybe_swap)

fibUp AGDA
* https://parfunc.slack.com/archives/C0BAZ6ZB8/p1498339139429555




\begin{figure}[t!]
\centering
\begin{minipage}[t]{0.30\textwidth}
\begin{codebox}
ex1 :: Nat -> Nat
ex1 x =
  let y =
    let t = x
    in
      dec t
  in
    inc y
$\wwedge$
\end{codebox}
\end{minipage}
\hspace{0.1in}
\begin{minipage}[t]{0.45\textwidth}
\begin{codebox}
ex2 :: Nat -> Nat
ex2 x =
  let ys = let n  = dec x
               p  = inc x
               xs = n : []
           in p : xs
      y  = last ys
  in
    inc y

\end{codebox}
\end{minipage}
\caption{Examples: (\bkw{ex1}) Local variables, (\bkw{ex2}) Collections}
\label{fig:ex1}
\label{fig:ex2}
\end{figure}


type Fusion = f:_ -> g:_ -> xs:_ ->
                {map (f . g) xs = (map f . map g) xs}

map_fusion f g []     = ()
map_fusion f g (x:xs) = map_fusion f g xs

map_fusion f g []
  =  (map f . map g) []
  =.  map f (map g [])
  =.  map f []
  =.  map (f . g) []
  **  QED
map_fusion f g (C x xs)
  =  map (f . g) (x : xs)
  =. ((f . g) x) : (map (f . g) xs)
  =. ((f . g) x) : (((map f) . (map g)) xs)
     ? map_fusion f g xs
  =. ((f . g) x) : (map f (map g xs))
  =. (f (g x)) : (map f (map g xs))
  =. map f ((g x) : (map g xs))
  =. map f ((g x) : (map g xs))
  =. map f (map g (x : xs))
  =. map f ((map g) (x : xs))
  =. ((map f) . (map g)) (x : xs)
  ** QED

type Assoc = xs:_ -> ys:_ -> zs:_ ->
              { xs ++ (ys ++ zs) = (xs ++ ys) ++ zs}

app_assoc              :: Assoc
app_assoc []     ys zs = ()
app_assoc (x:xs) ys zs = app_assoc xs ys zs


append_assoc [] ys zs
  = ([] ++ ys) ++ zs
  =. ys ++ zs
  =. [] ++ (ys ++ zs)
  ** QED

append_assoc (x:xs) ys zs
  = ((x : xs) ++  ys)  ++ zs
  =. (x : (xs ++  ys)) ++ zs
  =.  x :((xs ++  ys)  ++ zs)
  =.  x :((xs ++ (ys   ++ zs)) ? append_assoc xs ys zs
  =. (x : xs) ++ (ys   ++ zs)
  ** QED


append_assoc [] ys zs
  = ([] ++ ys) ++ zs
  =. ys ++ zs
  =. [] ++ (ys ++ zs)
  ** QED

append_assoc (x:xs) ys zs
  = ((x : xs) ++  ys)  ++ zs
  =. (x : (xs ++  ys)) ++ zs
  =.  x :((xs ++  ys)  ++ zs)
  =.  x :((xs ++ (ys   ++ zs)) ? append_assoc xs ys zs
  =. (x : xs) ++ (ys   ++ zs)
  ** QED


\mypara{Reflection vs. Axiomatization}
%
An alternative \emph{axiomatic} approach,
used by Dafny~\citep{dafny} and
\fstar~\citep{fstar},
is to encode @fib@ using a universally
quantified SMT formula (or axiom):
$\forall n.\ \fibdef\ (\fib\ n)\ n$.
%
Axiomatization offers greater automation than
reflection. Unlike \toolname, \fstar
%and \fstar
will verify the following by
\emph{automatically instantiating} the above
axiom at @2@, @1@ and @0@:
%
\begin{code}
   val axPf_fib2: unit -> Tot (v:unit {fib 2 = 1})
   let axPf_fib2 _ =  ()
\end{code}
%
The automation offered by axioms is a bit of a
devil's bargain, as axioms render checking of
the VCs \emph{undecidable}.
%
In practice, automatic axiom instantation can
easily lead to infinite ``matching loops''.
%
For example, the existence of a term \fib{n} in a VC
can trigger the above axiom, which may then produce
the terms \fib{(n-1)} and \fib{(n-2)}, which may then
recursively give rise to further instantiations
\emph{ad infinitum}.
%
To prevent matching loops an expert must carefully
craft ``triggers'' and provide a ``fuel''
parameter~\citep{Amin2014ComputingWA} that can be
used to restrict the numbers of the SMT unfoldings,
which ensure termination, but can cause the axiom
to not be instantiated at the right places.
%
As an example, \fstar's default fuel is not adequate
to compute the value of @fib@ at @15@,
rendering the following property unpredictably unsafe.
%
\begin{code}
   val axPf_fib2: unit -> Tot (v:unit {fib 15 = 610})
   let axPf_fib2 _ =  ()
\end{code}
%
In short, per the authors of Dafny, the
undecidability of the VC checking and its
attendant heuristics makes verification
unpredictable~\citep{Leino16}.




\subsection{Proof By Evaluation}

To simplify trivial proof by rewriting (like @fib 2 == 1@) while preserving
predictable behavior we introduce \pbename (Proof By Evaluation).
%
\pbename symbolically evaluates the subexpressions of the proof
obligations based on the definitions of the
reflected functions.
%
For instance the proof obligation
@fib 2 == 1@ fires the unfolding of $\fib{2}$
which in turn fires the unfoldings of
both $\fib{1}$ and $\fib{0}$.
%
All the three unfoldings are collected in an assumption
environment under which the proof obligation is checked
strengthening the initial obligation to the
following SMT query
$$
(\fib{2} = \fib{1} + \fib{0}
\land
\fib{1} = 1
\land
\fib{0} = 0
) \Rightarrow \fib{2} = 1
$$
%
Thus the SMT can prove @fib 2 == 1@
while treating @fib@ as an uninterpreted function.


\mypara{Inductive Proofs}
\pbename only unfolds reflected functions when it
can be statically decided exactly which branch the function unfolds to.
%
%
For example \pbename cannot automatically prove inductive properties like
that @fib@ is (locally) increasing: for all $n$, $\fib{n} \leq \fib{(n+1)}$.
%
\pbename will not automatically unfold \fib{n} since it is not
determined which branch @fib n@ will take, for all @n@.
%
The programmer can use standard analysis on the value of @n@
to appropriately fire \pbename unfoldings and thus construct inductive proofs.
%
\begin{mcode}
    fibUp :: n:Nat -> { !fib! n <= !fib! (n+1) }
    fibUp n | n == 0
            -- UNFOLDINGS
            -- fib 0 == 1, fib 1 == 1
            =  trivial
            | n == 1
            -- UNFOLDINGS
            -- fib 2 == fib 1 + fib 0, fib 1 == 1, fib 0 == 0
            = trivial
            | otherwise
            -- UNFOLDINGS
            -- fib (n+1) == fib n + fib (n-1)
            -- fib n == fib (n-1) + fib (n-2)
            =  fibUp (n-1) &&& fibUp (n-2)
\end{mcode} %$

\mypara{Case Splitting and Induction}
%
The proof @fibUp@ works by induction on @n@.
%
In the \emph{base} cases @0@ and @1@, we simply assert
the relevant unfoldings. These are verified as \pbename automatically
unfolds the
reflected @fib@ definition at those inputs.
%
The derived VCs are (automatically) proved
as the SMT solver concludes $0 < 1$ and $1 + 0 \leq 1$
respectively.
%
In the \emph{inductive} case, @fib n@ is automatically unfolded
to  @fib (n-1) + fib (n-2)@, which, because of the
induction hypothesis (applied by invoking @fibUp@
at @n-1@ and @n-2@) and the SMT solver's arithmetic
reasoning, completes the proof.

The connective operator @&&&@ is operationally the constant
function that moreover combines in the logic the information
form both its input arguments.

\begin{mcode}
    x &&& _ = x
\end{mcode}

\mypara{Properties of \pbename}
In \S~\ref{sec:pbe} we formalize
the automatically generated unfoldings
via \pbename and discuss three main
properties of our formalization
\begin{itemize}
\item \textbf{Completeness}
\pbename will actually unfold all
the reflected functions that can
be unfolded to exactly one branch.
\item \textbf{Termination}
Assuming that all the unfolded functions terminated,
\pbename will generate finitely many unfoldings.
\end{itemize}

As a result of the above two properties \pbename is more predictable
than fuel-based axiom instantiation approach that is traditionally
used by SMT-verifiers, since \pbename is automatically adjusting the
number of instantiations (\ie fuel) based on the context.
%
Automatic fuel adjustment though
may introduce a different source of
unpredictability: verification
slow-downs.
%
For instance verification of
@fib 16 == 987@ is unpredictably
unsafe in the fuel-based \fstar
and unpredictably slow in our
\pbename-based approach.
%
Yet, compared with axiomatization
our notion of unpredictability is
\textbf{local}: \pbename is separately
strengthening proof obligations thus
avoiding the ``butterfly effect''~\cite{Leino16}
where axioms introduced an one point
of the code may affect verification
in different places.

\mypara{Reflection and Refinement Types}
\NV{This seems out of place now. }
The above examples illustrate how refinement types
critically allow us to restrict the domain of both
reflected functions and proof terms in order to
enable sound theorem proving.
%
\begin{itemize}
  \item \emphbf{Totality}
  First, by refining @Int@ to @Nat@ we restrict
  the domain of Haskell's partial function @fib@
  which lets us use @!fib!@ as a total function in
  the logic. Specifically, refinement type checking
  ensures soundness by verifying that @fib@ is never
  called on negative numbers.
%
\item \emphbf{Induction}
  Second, the refinement type @Nat@
  in the argument of @fibUp@ constrains @fibUp@'s
  domain, turning @fibUp@ into a total proof that
  encodes induction on Natural numbers, while
  reasoning about Haskell's runtime integers.
\end{itemize}
%
Thus, refinement types let us encode proofs that
only hold on refined types, like @Nat@ while using
Haskell's Integers, instead of forcing to switch
to a user or library defined @Data.Natural@ Haskell
type.
%
As a less elegant alternative, usually followed
in dependently typed languages, @fib@ could be
defined as a function returning a @Maybe@ -- \ie
returning @Nothing@ on negative inputs.
%
\VC{I'm not sure if this claim could be made for dependently typed languages,
  for example, in Agda, Nat is compiled to Haskell Integer; in Coq, nat is
  compiled to OCaml int}
%
Then the code and proofs would be littered with cases analyses
required to extract @fib@'s result.
%
Another alternative is \dafny's approach
that uses code level assumptions and assertions
to restrict function's domain and range.
%
While these assumptions and assertions can
be semi-automatically discharged using the
SMT solver, they impose programmer overheads
as they cannot be inferred unlike refinement
types that are inferred using
the abstract interpretation framework of
liquid typing~\cite{LiquidPLDI08}.


\subsection{Higher Order Reasoning}

Refinement reflection allows terminating functions to appear in the refinements
but to preserve decidability refinements cannot contain quantifiers.
%
Next we present how our refinement types can be used to
encode quantified properties using lambda abstraction and dependent pairs
for forall and exists respectively,
thus encoding arbitrary higher order properties.

\mypara{Encoding foralls}
%
Refinements smoothly accomodate universal quantification as lambda abstraction.
%
For example, lets prove that every locally increasing
function is monotonic, \ie
if @f z <= f (z+1)@ for all @z@,
then @f x <= f y@ for all @x < y@.
%
\begin{mcode}
  fMono :: f:(Nat -> Int)
        -> fUp:(z:Nat -> {f z <= f (z+1)})
        -> x:Nat
        -> y:{x < y}
        -> {f x <= f y} / [y]
  fMono f inc x y
    | x + 1 == y = fUp x
    | x + 1 < y  = fMono f fUp x (y-1) &&& fUp (y-1)
\end{mcode}
%
We prove the theorem by induction
on @y@, which is specified by the
annotation @/ [y]@ which states
that @y@ is a well-founded
termination metric that decreases
at each recursive call~\citep{Vazou14}.
%
% All reflected functions are proved terminating.
% When the annotation metric is not explicit Liquid Haskell
% successfully uses heuristics to automatically prove termination.
%
If @x+1 == y@, then we use @fUp x@.
%
Otherwise, @x+1 < y@, and we use
the induction hypothesis \ie apply
@fMono@ at @y-1@, after which
transitivity of the less-than
ordering finishes the proof.
%
We can use the general @fMono@
theorem to prove that @fib@
increases monotonically:
%
\begin{code}
  fibMono :: n:Nat -> m:{n<m} -> {fib n <= fib m}
  fibMono = fMono fib fibUp
\end{code}

\mypara{Encoding existentials}
Refinements encode existential quantification as dependent pairs.
%
For example, we can prove that is there exists an @m@ greater than @n@
then there exists an upper bound to @fib n@.

\begin{mcode}
  fibBound :: n:Nat -> (m::Int, {n < m}) -> (m'::Int, {fib n <= m' }) @-}
  fibBound n (m, pf) = (fib m, fibMono n m)
\end{mcode}
\VC{$m::Int$ or $m:Int$?}
We use the notation @(m::t1, t2)@ to encode dependent pairs:
we bind the first element of the pair to the binder @m@,
then @m@ can appear in the refinements of the type @t2@ of the
second element.
%
To prove @fibBound@ we open the existential assumption to get a value
@m@ that is greater then @n@.
Then, @fib m@ bounds @fib n@ as proved by @fibMono n m@.

Even though quantifiers cannot appear inside the refinements
lambdas and dependent pairs allows us to quantify values over
the refinement properties.
%
One limitation of this encoding is that quantifiers cannot
exist inside logical connectives (like $\land$ and $\lor$).
%
In \S~\ref{sec:higher-order} we encode logical connectives using
data types.
For instance conjunction is encoded as a product and
disjunction as a union data type, thus naturally getting the
expressiveness power of arbitrary higher order logics like
Isabelle/HOL.


\subsection{Case Study: Peano Numerals}

Refinement reflection is not limited to programs
operating on integers. We conclude the overview
with a small library for Peano numerals, defined
via the following algebraic data type:
%
\begin{code}
  data Peano = Z | S Peano
\end{code}
%
We can @add@ two @Peano@ numbers via:
%
\begin{code}
  reflect add :: Peano -> Peano -> Peano
  add Z     m = m
  add (S n) m = S (add n m)
\end{code}
%
In \S~\ref{subsec:measures} we will describe
exactly how the reflection mechanism (illustrated
via @fibP@) is extended to account for ADTs like @Peano@.
%
Note that \toolname automatically checks
that @add@ is total~\citep{Vazou14}, which
lets us safely @reflect@ it into the
refinement logic.

\mypara{Add Zero to Left}
%
As an easy warm up, lets show that
adding zero to the left leaves the
number unchanged:
%
\begin{code}
  zeroL :: n:Peano -> { add Z n == n }
  zeroL n = trivial
\end{code}

\mypara{Add Zero to Right}
%
It is slightly more work to prove
that adding zero to the right also
leaves the number unchanged.
%
\begin{mcode}
  zeroR :: n:Peano -> { add n Z == n }
  zeroR Z     =  trivial

  zeroR (S n) =  zeroR n
\end{mcode}
%
The proof goes by induction, splitting cases on
whether the number is zero or non-zero. Consequently,
we pattern match on the parameter @n@, and furnish
separate proofs for each case.
%
In the ``zero" case, we simply unfold the definition
of @add@.
%
In the ``successor" case, after unfolding we (literally)
apply the induction hypothesis by using the because operator.
%
\toolname's termination and totality checker
verifies that we are in fact doing induction
properly, \ie the recursion in @zeroR@ is
well-founded~(\S~\ref{sec:types-reflection}).

\mypara{Commutativity}
%
Lets conclude by proving that @add@ is commutative:
%
\begin{mcode}
add_com :: a:_ -> b:_ -> {add a b = add b a}

add_com    Z  b =  zeroR b
add_com (S a) b =  add_com a b &&& sucR b a
\end{mcode}
%
using a lemma @sucR@
%
\begin{code}
  sucR :: n:Peano -> m:Peano ->
                {add n (S m) = S (add n m)}
  sucR = exercise_for_reader
\end{code}

%
% Again, note how case-splitting, induction, and
% lemma reuse are just pattern-matching, recursion and
% function application, via the lemma @sucR@:

Thus, refinement reflection lets us prove properties of Haskell programs
just by writing Haskell programs: lemmas are just functions, case-splitting
is just branching and pattern matching, and induction is just recursion.
%
Next, we formalize refinement reflection and describe
how to keep type checking decidable and predictable.
